Para confirmar que a eq. (5.22) é dimensionalmente correta, verifica-se que todos os termos têm a mesma dimensão física.

\textbf{Dimensões dos elementos do circuito RC:}
\begin{itemize}
\item Resistência $R$: $[R] = \Omega = \mathrm{V\cdot A^{-1}}$
\item Capacitância $C$: $[C] = \mathrm{F} = \mathrm{A\cdot s\cdot V^{-1}}$
\item Produto $RC$: $[RC] = \mathrm{s}$ (dimensão de tempo)
\item Tensão $V$: $[V] = \mathrm{V}$
\end{itemize}

\textbf{Verificação dimensional term a termo:}
A eq. (5.22) é do tipo $RC\,\dfrac{dV_C}{dt} + V_C = V_s$. Verificando cada parcela:
\[
\left[RC\,\frac{dV_C}{dt}\right] = [\mathrm{s}]\cdot\left[\frac{\mathrm{V}}{\mathrm{s}}\right] = \mathrm{V},
\]
\[
[V_C] = \mathrm{V},\qquad [V_s] = \mathrm{V}.
\]
Portanto, todos os termos têm dimensão de tensão (volts):
\[
[RC\,dV_C/dt]\;=\;[V_C]\;=\;[V_s]\;=\;\mathrm{V}.\quad\checkmark
\]
\textbf{Conclusão:} A eq. (5.22) é dimensionalmente correta e consistente com a análise de circuitos elétricos. 